\section{Conjuntos Numéricos}
%\linespacedtitle
\subsection{Números Naturais e Inteiros}
%\linespaced
Vejamos abaixo conjuntos dos naturais e inteiros, representados por $\mathbb{N}$ e $\mathbb{Z}$, respectivamente.\\

\text{\textbullet} Naturais: $\mathbb{N} \; = \; \{ 1,2,3,4, ... \} $ \\

\text{\textbullet} Inteiros: $ \mathbb{Z} \; = \; \{ ..., -3, -2, -1, 0, 1, 2, 3, ... \}$ \\
%\linespaced

\noindent
\textbf{Propriedades} \\

%\linespaced

\noindent
Consideremos $a$, $b$ e $c$ números inteiros, \\
\noindent
\emph{A1.} $ (a + b) + c = a + (b + c)$ \\
\emph{A2.} $ a + b = b + a $ \\
\emph{A3.} $ a + 0 = a $ \\
\emph{A4.} $ (ab)c = a(bc) $ \\
\emph{A5.} $ ab = ba $ \\
\emph{A6.} $ a \cdot 1 = a $ \\
%\linespaced

\noindent
Inteiros positivos $\mathbb{Z_+}$ e inteiros negativos $ \mathbb{Z_{-}}$,\\

\noindent
$ \mathbb{Z_+} \; = \; \{ 0, 1, 2, 3, 4, ... \} $ e $ \mathbb{Z_{-}} \; = \; \{ ..., -4, -3, -2, -1, 0 \} $ \\

\noindent
\underline{Simétrico}: $ a \in \mathbb{Z} \; : \; a + (- a) = 0$, com $(-a) \in \mathbb{Z} $, \\
$ a - b = a + (-b) $ \\

\noindent
\underline{Divisor}: $a|b \Rightarrow \{ \exists c \in \mathbb{Z} | \; ca=b \}$, onde $a$ é divisor de $b$. \\ \\

%\linespaced

\noindent
\subsection{Números Racionais}
%\linespacedtitle

\noindent
Representamos conjunto dos racionais com a letra $\mathbb{Q}$, que são escritos na forma de fração, como definido abaixo, \\

\noindent
$\mathbb{Q} = \{ a \in \mathbb{Z}, b \in \mathbb{Z^*}| \; \frac{a}{b} \in \mathbb{Q} \} $ \\

\noindent
\emph{A1.} $ \frac{a}{b} \; = \; \frac{c}{d}$  $\Leftrightarrow ad = bc $ \\
\noindent
\emph{A2.} $\frac{a}{b} + \frac{c}{d} \; = \; \frac{ ad + bc }{bd}$ \\
\noindent
\emph{A3.} $ \frac{a}{b} \cdot \frac{c}{d} \; = \; \frac{ac}{bd}$ \\


\noindent
$ \mathbb{Q^{'}} \; = \; \{ x = \frac{a}{d} \in \mathbb{Z} | \ b = 1, \; x, a, b \in \mathbb{Z} \} = \mathbb{Z}  $ \\
$ \Rightarrow \mathbb{Z} \subset \mathbb{Q} $ \\

A multiplicação e a adição dos racionais também obedecem as regras de corpo, associatividade, comutatividade, elemento neutro, distributividade, elemento inverso e elemento simétrico.\\

\noindent
\textbf{Propriedades} \\

\noindent
\emph{A1.} $ (\frac{a}{b} + \frac{c}{d}) + \frac{e}{f}  = $ $\frac{a}{b} + ( \frac{c}{d} + \frac{e}{f} )$ \\

\noindent
\emph{A2.} $ \frac{a}{b} + \frac{c}{d} \; = \; \frac{c}{d} + \frac{a}{b} $ \\

\noindent
\emph{A3.} $ \frac{a}{b} + 0 \; = \; \frac{a}{b}$ \\

\noindent
\emph{A4.} $ \frac{a}{b} + ( - \frac{a}{b} ) \; = \; 0$ \\

\noindent
\emph{M1.} $ (\frac{a}{b} \cdot \frac{c}{d} ) \cdot \frac{e}{f} \; = \; \frac{a}{b} \cdot ( \frac{c}{d} \cdot \frac{e}{f} )$ \\

\noindent
\emph{M2.} $ \frac{a}{b} \cdot \frac{c}{d} \; = \; \frac{c}{d} \cdot \frac{a}{b}$ \\

\noindent
\emph{M3.} $ \frac{a}{b} \cdot 1 \; = \; \frac{a}{b} $ \\

\noindent
\emph{M4.} $ \frac{a}{b} \cdot ( \frac{c}{d} + \frac{e}{f}) \; = \; \frac{a}{b} \cdot \frac{c}{d} + \frac{a}{b} \cdot \frac{e}{f}  $ \\

\noindent
\emph{M5.} $ \frac{a}{b} \cdot (\frac{a}{b})^{-1} \; = \; \frac{a}{b} \cdot \frac{b}{a} \; = \; 1$ 


\subsection{Números Irracionais}
%\linespacedtitle

\noindent
São números irracionais, representado por $\mathbb{I}$, números $x$ tal que $x \notin \mathbb{Q}$, ou seja não são racionais e pertencem ao conjunto dos reais.

$$ \mathbb{Q} \cap \mathbb{I} = \emptyset$$ 

%\linespacedEndSection

\subsection{Números Reais}
%\linespacedtitle
\noindent
O conjuntos dos reais é formado pela união dos racionais e dos irracionais, 
$ \mathbb{R} = \mathbb{Q} \cup \mathbb{I} $.\\
%\linespaced

\noindent
\text{\maltese} Vamos demonstrar que $\; \sqrt{2} \notin \mathbb{Q}$ 

\begin{proof}
	Supondo $ \frac{a}{b} = \sqrt{2} \Rightarrow a^2 = 2b^2 $ \\
	$ \Rightarrow a^2 $ é par $ \Rightarrow $ $a$ é par $ \Rightarrow a = 2m$, $ m \in \mathbb{Z} \Rightarrow $ \\
	$ \Rightarrow (2m)^2 = 2b^2 \Rightarrow b^2 = 2m^2 \Rightarrow b \text{ é par } \Rightarrow $ \\
	logo um absurdo (supomos $\frac{a}{b}$ irredutível)  pois $mdc(a,b) = 1$. \\
\end{proof}

%\linespaced
\noindent
Obtendo novos números irracionais,\\

\noindent
1) $ \sqrt{p}$, onde $p$ é primo e positivo;\\
2) $\alpha \cdot r$, $ \alpha + r$, $ \frac{\alpha}{r}$, onde $\alpha$ é irracional e $r$ é racional ($ \alpha = \sqrt{p} $).\\

São dois métodos para se obter números irracionais.
%\linespaced
\vspace{18pt}

% Retas dos reais
\begin{tikzpicture}[scale=1.2,domain=-1:3/2]
	\draw[very thin, color=gray] (-2,-0.1) grid (5,0.1);
	\draw [->] (-2,0) -- (5, 0) node[right] {$\mathbb{R} \; \text{(reta real)}$};
	\draw (-1.0, -0.7 ) node[above] {$-1$};
	\draw (0, -0.7 ) node[above] {$\frac{-1}{2}$};
	\draw (1, -0.7 ) node[above] {$0$};
	\draw (2, -0.7 ) node[above] {$\frac{1}{2}$};
	\draw (3, -0.7 ) node[above] {$1$};
	\draw (4, -0.7 ) node[above] {$\frac{3}{2}$};
	\draw (3.8, 0.25 ) node[above] {$\mathit{\sqrt{2}}$};
	\draw [-] (3.8, -0.075) -- (3.8, 0.075); 

\end{tikzpicture}

A reta real é ordenada, um número $x$ é maior que qualquer número à sua esquerda e menor que qualquer número à sua direita. \\

\noindent
\underline{Intervalos Numéricos}\\
%\linespacedtitle

\noindent
\textbullet \; Intervalo aberto: $( a, b) = \{ x \in \mathbb{R} | \ a < x < b \};$ \\
\textbullet \; Intervalo fechado: $ [a, b] = \{ x \in \mathbb{R} | \ a \leq x \leq b  \} $ \\
\textbullet \; Intervalo fechado à esquerda e aberto à direita: $ [a, b) = \{ x \in \mathbb{R} |  \ a \leq x < b \} $ \\
\textbullet \; Intervalo aberto à esquerda e fechado à direita:  $ (a, b] = \{ x \in \mathbb{R} | \ a < x \leq b \} $